% Limits of the Approach
\section{Limits of the Approach}

\subsection{Inherent Limitations}

\subsubsection{Corpus Size and Quality}
The Indus corpus is:
\begin{itemize}
    \item Small (approximately 4,000 inscriptions)
    \item Short (median length 4--5 signs)
    \item Fragmentary (many damaged or partially preserved)
    \item Contextually limited (primarily from administrative and commercial contexts)
\end{itemize}

These constraints limit statistical confidence and restrict generalizations.

\subsubsection{Lack of Bilingual Texts}
Without Rosetta Stone-equivalent bilingual inscriptions, certain questions remain unanswerable:
\begin{itemize}
    \item Phonetic values of signs
    \item Semantic content beyond operational function
    \item Language affiliation
    \item Reading order (though strong evidence suggests left-to-right)
\end{itemize}

\subsection{What We Cannot Know}

This methodology cannot determine:
\begin{enumerate}
    \item The pronunciation or phonetic values of signs
    \item The specific natural language(s) represented
    \item Semantic meanings beyond operational categories
    \item Individual names, places, or historical events
    \item Cultural or religious significance of specific symbols
\end{enumerate}

\subsection{What We CAN Know}

Despite these limitations, we can rigorously demonstrate:
\begin{enumerate}
    \item The mathematical structure of the operational system
    \item The logical constraints on valid sequences
    \item The functional categories of signs
    \item The computational power of the system (equivalent to DFA)
    \item The administrative and validational purposes of inscriptions
\end{enumerate}

\subsection{Epistemological Position}

We distinguish between:
\begin{itemize}
    \item \textbf{Demonstrated facts}: Claims supported by distributional evidence and formal proof
    \item \textbf{Working hypotheses}: Plausible interpretations consistent with evidence but not uniquely determined
    \item \textbf{Speculation}: Interesting possibilities beyond current evidence
\end{itemize}

This work restricts itself primarily to the first category, with clearly marked excursions into the second.
